\section{Full component ablation study}\label{sec:ablation}

All 432 stages are complete: 54 verified full references and 378 new component evaluations. Figure~\ref{fig:ablations} displays ablation-minus-full effects. A positive interval indicates degradation after removal; a negative interval indicates an improvement. Thirteen exploratory intervals are wholly positive, four wholly negative and 25 include zero. Table~\ref{tab:ablations} identifies which components fall into these categories.

Stale blending has the clearest repeated benefit. Removing it worsens score in METR-LA C3, both PEMS-BAY conditions and both SUMO conditions. The largest effect occurs in PEMS-BAY C4: removal increases score by $4.220\,\ms$, with 95\% interval [3.760, 4.699]\,$\ms$, while coverage falls to 0.549. Spatial pooling is beneficial in METR-LA C3 and PEMS-BAY C3/C4, but its remaining intervals include zero.

The adverse findings matter equally. Removing age weighting improves score in METR-LA C4 and PEMS-BAY C3/C4; removing context matching improves PEMS-BAY C4. Consequently, the ablations do not support a claim that every component is necessary or uniformly beneficial. Arrival-time freshness is worse in METR-LA C4 and PEMS-BAY C3, but the remaining effects are small or uncertain. The simulated diagnostic does not establish a large general backlog penalty. Exploratory intervals are not used to modify the tested method.


Tables~\ref{tab:ablation-metr}--\ref{tab:ablation-sumo} provide all forty-two numerical removal effects, exploratory intervals, coverage and widths. Table~\ref{tab:ablation-metr} shows that removing age weights improves METR-LA C4 score, while other small effects remain inconclusive. Table~\ref{tab:ablation-pems} shows age-weight improvements in both PEMS-BAY cells and a context-removal improvement in C4. The same table documents the severe stale-blend failure in PEMS-BAY C4. Table~\ref{tab:ablation-sumo} shows a useful stale blend in both SUMO conditions, with several remaining intervals spanning zero. These findings are more informative than treating every component as necessary.

Figure~\ref{fig:ablations} uses separate panel scales to preserve small effects while explicitly labelling the off-scale PEMS-BAY C4 stale-blend difference. Its zero line corresponds to no paired score change. A beneficial removal is evidence against claiming that the removed component uniformly improves quality. Conversely, an inconclusive interval does not prove redundancy. Dependence between components, limited episode/day replication and multiple exploratory comparisons prevent a simple causal ranking of design choices. The locked full method is not retrospectively retuned to maximise these test-set results.

\subsection{Descriptive archive-cap sensitivity}

Eighteen timed replays compare archive caps 64, 128 and 256 with identical C3 inputs and fixed model/fault seeds 11/101. Table~\ref{tab:runtime} reports wall-clock medians, observed ranges and pooled quality. Timing includes CPU calibration, loading and output serialisation, excluding training and prediction. The cap-256 setting exactly reproduces the archived same-run quality. Smaller caps have slightly higher score in both fixed runs, but timing ranges overlap. Host workload and filesystem cache are not isolated, so these measurements do not establish a reliable speed advantage, a cross-method runtime ranking or a deployment throughput guarantee.


Figure~\ref{fig:compute} complements Table~\ref{tab:runtime}: horizontal ranges show three observed replay times at each cap, and colour/labels identify the cap rather than another forecasting method. The score axis uses the same-run pooled quality estimand. No error bar is a runtime confidence interval. This figure is a reproducibility and resource sensitivity diagnostic, not a speed competition among calibrators or evidence of service-level latency.
